
\subsubsection*{5.2.2}
\phantomsection
\label{sga3:VIIB:5.2.2}

Suppose first that \(B\) has height at most \(1\); that is,
\[
  x^p=0
  \qquad\text{for every }x\in I.
\]
By \SGARefLink{sga3:VIIB:5.1.4}{5.1.4}, the graded algebra
\[
  \operatorname{Gr}_I(B)
\]
associated with the filtration
\[
  I\supset I^2\supset I^3\supset\cdots
\]
of \(B\) is endowed with a cogroup structure in the category
\(\mathbf{Alpg}/k\). In other words, the profinite algebra
\[
  C=\prod_{n\geq0}\operatorname{Gr}_I(B)_n
\]
is the affine algebra of a formal \(k\)-group \(N\).
Clearly,
\[
  \omega_{N/k}=I/I^2,
\]
and \(N\) is infinitesimal of height at most \(1\).

By \SGARefLink{sga3:viib:subsection:4-4}{4.4}, the identity map of \(\omega_{N/k}\) therefore induces
an isomorphism
\[
  B'
  =
  k[[\omega_{N/k}]]
  \big/
  \bigl((x^p)\bigr)_{x\in\omega_{N/k}}
  \xrightarrow{\ \sim\ }
  C.
\]
In particular, the map \(\phi\) of \SGARefLink{sga3:VIIB:5.2.1:section-proposition}{5.2.1} induces an isomorphism
between the associated graded algebras of \(B'\) and \(B\), when
both are filtered by the powers of their augmentation ideals.
It follows from \textup{[CA]}, \(\S\)V.7, Lemma~1
\textup{(see also \SGARefLink{sga3:VIIB:5.1.5}{5.1.5})} that \(\phi\) is an isomorphism.
